\documentclass[11pt]{article}

\usepackage[a4paper,margin=1in]{geometry}
\usepackage{amsmath,amssymb,amsthm,mathtools}
\usepackage{array}
\usepackage{booktabs}
\usepackage{enumitem}
\usepackage{tikz}
\usepackage{physics}
\usepackage{needspace}
\usepackage{xcolor}
\definecolor{courseblue}{RGB}{0,90,150}
\usetikzlibrary{arrows.meta,positioning}

\setlength{\parindent}{0pt}
\setlength{\parskip}{0.75em}

\newcounter{problem}
\newenvironment{problem}{%
    \refstepcounter{problem}%
    \Needspace{7\baselineskip}%
    \par\vspace{0.7em}%
    \noindent\textcolor{courseblue}{\rule{\textwidth}{0.4pt}}%
    \par\vspace{0.5em}%
    \begin{list}{}{%
        \setlength{\leftmargin}{2.3em}%
        \setlength{\labelwidth}{1.8em}%
        \setlength{\labelsep}{0.5em}%
        \setlength{\topsep}{0pt}%
        \setlength{\partopsep}{0pt}%
    }%
    \item[{[\theproblem]}]
}{%
    \end{list}
    \par\vspace{0.5em}
}

\newenvironment{parts}{%
    \begin{enumerate}[label={[\alph*]},leftmargin=2.2em,itemsep=0.65em,topsep=0.65em]
}{%
    \end{enumerate}
}

% Type each solution between \begin{answer} and \end{answer}.
\newenvironment{answer}{%
    \Needspace{4\baselineskip}%
    \par\vspace{0.35em}
    \noindent\textbf{Answer:}\par
}{%
    \par\vspace{0.25em}
}



\title{CO250 Assignment 5}
\author{}
\date{October 19--November 2}

\begin{document}
\maketitle

\begin{problem}
    Consider the linear program
    \begin{equation*}
    \begin{aligned}
        \max\quad & x_4 \\
        \text{s.t.}\quad & x_1-x_3=1, \\
        & x_2-x_3=-2, \\
        & x_1-x_2-x_3-x_4=1, \\
        & x_1,x_2,x_3,x_4\geq 0.
    \end{aligned}
    \tag{P}
    \end{equation*}
    Equivalently, $P$ has the form
    \[
        \max\{c^Tx:Ax=b,\ x\geq\mathbf{0}\},
    \]
    where
    \[
        A=\begin{bmatrix}1&0&-1&0\\0&1&-1&0\\1&-1&-1&-1\end{bmatrix},
        \qquad b=\begin{bmatrix}1\\-2\\1\end{bmatrix},
        \qquad c=\begin{bmatrix}0\\0\\0\\1\end{bmatrix}.
    \]
    \begin{parts}
        \Needspace{7\baselineskip}\item Write down the auxiliary LP for (P).
        \begin{answer}

        \end{answer}
        \Needspace{10\baselineskip}\item Run the first phase of Simplex method for solving (P).
        Apply Bland's rule whenever there are multiple choices for an entering or leaving variable.
        (The Bland's rule says that if there are multiple variables that could enter or leave the basis,
        always choose the one with the smallest index.)
        \begin{answer}

        \end{answer}
        \Needspace{8\baselineskip}\item Explain how to obtain a feasible basis for the second phase of the simplex method for solving (P).
        \begin{answer}

        \end{answer}
        \Needspace{10\baselineskip}\item There are other ways of constructing an auxiliary LP.
        Find an auxiliary LP for (P) that has different objective function and/or different constraints from the one in part (a).
        Briefly explain how to determine the feasibility of (P) by solving the new auxiliary LP.
        \begin{answer}

        \end{answer}
    \end{parts}
\end{problem}

\begin{problem}
    In this question we study a restricted family of LP's whose feasible region is given only by
    linear equations, without the non-negativity constraint we usually have in SEF.
    For a matrix $A\in\mathbb{R}^{m\times n}$, the \emph{kernel} $\ker(A)$ of $A$ is defined by
    $\ker(A)=\{x\in\mathbb{R}^n:Ax=0\}$. Consider the LP
    \begin{equation}
    \begin{aligned}
        \max\quad & c^Tx \\
        \text{subject to}\quad &
        \begin{pmatrix}2&-3&-4\\1&5&11\end{pmatrix}x
        =\begin{pmatrix}-5\\17\end{pmatrix},
    \end{aligned}
    \label{eq:restricted-lp}
    \end{equation}
    which features in assignment 3, question 1b, with $c=(-1,-1,-1)^T$.
    % The source's reference to assignment 3, question 1b is preserved verbatim.
    \begin{parts}
        \Needspace{12\baselineskip}\item Show that
        \[
            \begin{pmatrix}2&-3&-4\\1&5&11\end{pmatrix}d=0
            \text{ if and only if }
            d=t\begin{pmatrix}1\\2\\-1\end{pmatrix}
            \text{ for some }t\in\mathbb{R}.
        \]
        \begin{answer}

        \end{answer}
        \Needspace{8\baselineskip}\item Show that if $c\in\mathbb{R}^3$ is such that $c^T(1,2,-1)\neq 0$, then the LP \eqref{eq:restricted-lp} is unbounded.
        \begin{answer}

        \end{answer}
        \Needspace{8\baselineskip}\item Show that if the LP \eqref{eq:restricted-lp} is unbounded, then $c^T(1,2,-1)\neq 0$.
        \begin{answer}

        \end{answer}
        \Needspace{9\baselineskip}\item Consider an LP of the form $\max\{c^Tx:Ax=b\}$, which we assume to be feasible.
        Prove that LP has an optimal solution if and only if $c$ is orthogonal to every vector in $\ker(A)$.
        \begin{answer}

        \end{answer}
    \end{parts}
\end{problem}

\Needspace{16\baselineskip}
\begin{problem}
    Consider the hyperplanes in $\mathbb{R}^3$
    \[
        H_1:\ x+y+z=2,\qquad H_2:\ x-y+z=1.
    \]
    \begin{parts}
        \Needspace{7\baselineskip}\item Find a normal vector to each hyperplane.
        \begin{answer}

        \end{answer}
        \Needspace{7\baselineskip}\item Determine the acute angle between the two hyperplanes. Show your calculation.
        \begin{answer}

        \end{answer}
        \Needspace{14\baselineskip}\item Consider the polyhedron
        \[
            P=\left\{(x,y,z)\in\mathbb{R}^3:
            \begin{aligned}
                x+y+z&\leq 2,\\
                x-y+z&\leq 1,\\
                x,y,z&\geq 0
            \end{aligned}\right\}.
        \]
        Explain why $P$ is a polyhedron, and identify the half-spaces defining $P$.
        \begin{answer}

        \end{answer}
        \Needspace{8\baselineskip}\item Replace the inequality $x-y+z\leq 1$ by $x-y+z>1$.
        Does $P$ remain a polyhedron? Justify your answer.
        \begin{answer}

        \end{answer}
    \end{parts}
\end{problem}

\Needspace{23\baselineskip}
\begin{problem}
    For points $x^1,x^2$ in $\mathbb{R}^n$ we define
    \[
        L_{x^1,x^2}=\{\lambda x^1+(1-\lambda)x^2:0\leq\lambda\leq 1\}
    \]
    to be the \emph{line segment} connecting $x^1$ and $x^2$.
    We then say that a set $S\subseteq\mathbb{R}^n$ is a \emph{convex} if for every pair of points
    $x^1\in S$ and $x^2\in S$ the line segment $L_{x^1,x^2}$ is fully contained in $S$.

    Let $x_1,x_2,x_3\in\mathbb{R}^n$ be three given vectors. Show that the set
    \[
        \Delta=\left\{\sum_{i=1}^3\mu_i x_i:
        \mu_i\geq 0\text{ for all }i,\sum_{i=1}^3\mu_i=1\right\}
    \]
    is convex. Use \underline{only} the above definition of convexity!
    \begin{answer}

    \end{answer}
\end{problem}

\end{document}
